\section[B. AIMC实验室]{\myboxii{      \BF{\TT{B}}} AIMC实验室}

\textit{运行时间限制}   X second

\textit{内存使用限制} X Gigabyte

\textit{读入文件} \qquad 标准输入

\textit{输出文件} \qquad 标准输出


\begin{quote}
	\textit{街灯璀璨，可我日渐透明。}
\end{quote}

\subsection{题目描述}

你所在的 AIMC 实验室研究出了一种全新的三维量子矩阵！

给定一个大小为 $N \times M \times H$ 的三维空间网格，坐标范围为 $1 \le x \le N$, $1 \le y \le M$, $1 \le z \le H$。我们用 $u$ 代表坐标为 $(x_u, y_u, z_u)$ 的量子。

对于网格中的每一个量子，你可以将其初始化为两种基态之一：$|0\rangle$ 态或 $|1\rangle$ 态。

对于任意两个量子 $u$ 和 $v$，定义它们之间的\BF{平方空间距离}为：
$$
d^2(u, v) = (x_u - x_v)^2 + (y_u - y_v)^2 + (z_u - z_v)^2
$$

系统存在 $8$ 个大小同为 $N \times M \times H$ 的能量矩阵 $A_1, A_2, \ldots, A_8$。系统总能量计算规则如下：

对于网格中的任意量子 $u$：

\begin{itemize}
	\item 若 $u$ 为 $|0\rangle$ 态，总能量增加 $A_1(u)$ ；
	\item 若 $u$ 为 $|1\rangle$ 态，总能量增加 $A_2(u)$ 。
\end{itemize}

定义 $S_1(u)$ 为网格内所有满足 $d^2(u, v) = 1$ 的量子 $v$ 的集合：

\begin{itemize}
	\item 若 $u$ 为 $|0\rangle$ 态，且 $S_1(u)$ 中所有量子均为 $|1\rangle$ 态，总能量增加 $A_3(u)$ ；
	\item 若 $u$ 为 $|1\rangle$ 态，且 $S_1(u)$ 中所有量子均为 $|0\rangle$ 态，总能量增加 $A_4(u)$ 。
\end{itemize}

定义 $S_2(u)$ 为网格内所有满足 $d^2(u, v) = 2$ 的量子 $v$ 的集合：

\begin{itemize}
	\item 若 $u$ 为 $|0\rangle$ 态，且 $S_2(u)$ 中所有量子均为 $|0\rangle$ 态，总能量增加 $A_5(u)$ ；
	\item 若 $u$ 为 $|1\rangle$ 态，且 $S_2(u)$ 中所有量子均为 $|1\rangle$ 态，总能量增加 $A_6(u)$ 。
\end{itemize}

定义 $S_3(u)$ 为网格内所有满足 $d^2(u, v) = 3$ 的量子 $v$ 的集合：

\begin{itemize}
	\item 若 $u$ 为 $|0\rangle$ 态，且 $S_3(u)$ 中所有量子均为 $|1\rangle$ 态，总能量增加 $A_7(u)$ ；
	\item 若 $u$ 为 $|1\rangle$ 态，且 $S_3(u)$ 中所有量子均为 $|0\rangle$ 态，总能量增加 $A_8(u)$ 。
\end{itemize}

作为AIMC实验室的首席科学家，请你为网格中的每个量子指定一个状态，使得系统总能量最大化，并输出这个最大值。

\subsection{输入格式}

第一行输入三个整数 $N$, $M$, $H$ ($1 \le N, M, H \le 20$)，表示三维网格的大小。

接下来依次给出 8 个能量观测矩阵 $A_1$ 到 $A_8$ 的数据。

对于每个矩阵，包含 $N$ 个区块。每个区块包含 $M$ 行，每行 $H$ 个整数。区块之间用空行隔开。

第 $x$ 个区块的第 $y$ 行第 $z$ 个数即代表对应坐标 $(x, y, z)$ 的能量值，所有矩阵中的能量值 $a_i \in [0, 10^9]$。

\subsection{输出格式}

输出一行一个整数，表示系统可以获得的最大总能量值。

\subsection{样例 1}

\begin{minipage}{\linewidth}
	\begin{minipage}[t]{0.49\linewidth}
		\begin{lstlisting}[title=\texttt{standard input}]
1 1 1
3

5

100

200

1000

2000

50

60
\end{lstlisting}
	\end{minipage}
	\begin{minipage}[t]{0.49\linewidth}
		\begin{lstlisting}[title=\texttt{standard output}]
2265
\end{lstlisting}
	\end{minipage}
\end{minipage}


\subsection{样例 1 解释}

整个网格仅有一个量子 $(1, 1, 1)$。

将该量子初始化为 $|1\rangle$ 时获得 $A_2 + A_4 + A_6 + A_8 = 5 + 200 + 2000 + 60 = 2265$，为最优。

\subsection{样例 2}

\begin{minipage}{\linewidth}
	\begin{minipage}[t]{0.49\linewidth}
		\begin{lstlisting}[title=\texttt{standard input}]
2 2 2
1 1
1 1

1 1
1 1

0 0
0 0

0 0
0 0

0 0
0 0

0 0
0 0

0 0
0 0

0 0
0 0

10 10
10 10

10 10
10 10

0 0
0 0

0 0
0 0

0 0
0 0

0 0
0 0

0 0
0 0

0 0
0 0
\end{lstlisting}
	\end{minipage}
	\begin{minipage}[t]{0.49\linewidth}
		\begin{lstlisting}[title=\texttt{standard output}]
88
\end{lstlisting}
	\end{minipage}
\end{minipage}


\subsection{样例 2 解释}

将 $2 \times 2 \times 2$ 中所有 $8$ 个量子统一初始化为 $|0\rangle$。

\begin{itemize}
	\item 每个量子从 $A_1$ 各获得 $1$ 的能量，合计 $8$。
	\item 每个量子距离为 $2$ 的邻居也全部为 $|0\rangle$，从 $A_5$ 各获得 $10$，合计 $80$。
\end{itemize}

总能量 $8 + 80 = 88$。可以证明这是最大值。
